\section{Open Questions, Ambiguities, and Likely Failure Points}

\subsection{Concepts That Are Useful but Still Fuzzy}

Several ideas in the knowledge base are undeniably valuable, but still too fuzzy to treat as fully specified rules:
\begin{itemize}
  \item ``real'' versus ``fake'' density;
  \item ``clean'' versus ``unclean'' trade examples;
  \item whether a robot is still active or already done;
  \item how much post-break delay is still acceptable before calling failure;
  \item how many tests weaken a level enough to invalidate bounce logic;
  \item what counts as meaningful support behind a visible order.
\end{itemize}

These should remain explicitly marked as ambiguous in the research code and documentation.

\subsection{Potential Contradictions Inside the Source Worldview}

The materials are mostly coherent, but several tensions exist:
\begin{itemize}
  \item visible size is treated as critical, yet also treated as easy to fake;
  \item heavy aggression can confirm continuation, yet also mark exhaustion;
  \item repeated tests can strengthen a breakout thesis, yet too many tests can degrade level quality;
  \item strong trends invite breakout participation, yet stretched movement invites spring reversals.
\end{itemize}

These are not fatal contradictions.
They mean the state space is conditional, not binary.

\subsection{Ideas That Need Empirical Testing Most Urgently}

\begin{enumerate}
  \item Whether visible-liquidity bounce logic survives modern crypto and modern Russian-market conditions.
  \item Whether order-book-only archive data is sufficient to say anything useful about breakout versus fakeout transitions.
  \item Whether simple day-specific adaptation materially improves detector quality.
  \item Whether robot-like patterns can be detected robustly enough to trade systematically rather than anecdotally.
  \item Whether cluster-memory alignment adds measurable signal once live logging is rich enough.
\end{enumerate}

\subsection{What May Be Outdated}

Some parts of the materials may be era-specific:
\begin{itemize}
  \item terminal-specific workflow assumptions;
  \item size thresholds calibrated to older MOEX liquidity conditions;
  \item participant behavior patterns that have changed;
  \item and infrastructure assumptions that matter differently in crypto.
\end{itemize}

The repo should preserve the ideas but not worship their original numeric details.

\subsection{Practical Conclusion}

The knowledge base is strongest when treated as a practitioner hypothesis library.
Its main value is not that it hands over a finished algorithm.
Its value is that it tells the repository where to look:
\begin{itemize}
  \item at defended liquidity,
  \item at compression and release,
  \item at the relation between displayed size and aggressive flow,
  \item at participant behavior patterns,
  \item and at day-specific adaptation.
\end{itemize}

The correct next step is therefore not ``optimize a final strategy''.
The correct next step is to keep building a platform that can:
\begin{itemize}
  \item log cleanly,
  \item replay honestly,
  \item detect simple patterns explicitly,
  \item and compare historical assumptions against live observation.
\end{itemize}
